\documentclass[11pt]{article}

\usepackage[margin=1in]{geometry}
\usepackage{graphicx}
\usepackage{amsmath}
\usepackage{amssymb}
\usepackage{authblk}
\usepackage[hidelinks]{hyperref}
\usepackage{natbib}
\usepackage{booktabs}
\usepackage{caption}
\captionsetup{font=small,labelfont=bf}

\renewcommand\Authfont{\normalsize}
\renewcommand\Affilfont{\small\itshape}

\title{\bfseries A heavy-tailed law of contextual memory in human language}

\author[1]{Elan Barenholtz}
\affil[1]{Center for the Future Mind, Florida Atlantic University}

\date{}

\begin{document}
\maketitle

\begin{abstract}
\noindent
Many human behaviors exhibit power-law temporal structure with no characteristic
scale. Whether language shares this signature has not been directly measured. We
use large language models as probabilistic probes to measure how much prior
context constrains the probability of upcoming tokens, isolated from
token-frequency effects through a shuffled-token baseline. The persistence
function $P(d)$ quantifies this order-specific influence at distance $d$ from the
current point of generation. Across ten corpora spanning six language families,
multiple genres, and written and spoken modalities, $P(d)$ decays approximately
as $1/d$, with mean slope $-1.04$ (SD $=0.15$) across nearly two orders of
magnitude, vanishes in random-token controls, replicates across independent probe
models, and is broadly distributed across prior sentences rather than concentrated
in anchors. Order- and content-driven contributions are distinguishable, with
content-driven integration carrying the long-range tail. The continuous scale-free
form is difficult to reconcile with architectural accounts that posit a sharp
division between maintained and retrieved material, and is consistent instead with
memory in language as a graded distance-dependent influence carried by the evolving
probability structure of the sequence itself.
\end{abstract}

\section{Introduction}

Human behavior across many domains exhibits lawful temporal structure across many
timescales. In neural activity \citep{beggs2003neuronal}, motor and cognitive
performance \citep{gilden2001cognitive}, and inter-event waiting times in human
dynamics \citep{barabasi2005origin}, distance-dependent dependencies decay as power
laws rather than as bounded or exponential functions, placing these processes among
the broad class of human behaviors with scale-free temporal organization
\citep{kello2010scaling}. Language is among the longer-range coherent sequential
behaviors humans produce, sustaining structure across thousands of words, but the
quantitative form of its temporal structure --- how the influence of prior material
scales with distance from the current point of generation --- has not been directly
measured. Existing characterizations of contextual integration in language,
including discourse models, situation models, and resonance-based accounts, operate
at the representational rather than the quantitative level, and conventional measures
of statistical dependence in text conflate sequential structure with vocabulary-level
regularities.

Here we develop a distance-resolved measurement of contextual influence in natural
language and apply it across ten corpora spanning six language families, multiple
genres, and written and spoken modalities. We use large language models as
probabilistic probes, treating their conditional probability estimates as a
calibrated readout of how preceding context constrains upcoming tokens. Because naive
perplexity reduction conflates sequential dependence with the probe's distributional
calibration, we subtract a shuffled-token baseline that absorbs the contribution of
token identity alone. The resulting persistence function $P(d)$ quantifies the
order-specific influence of prior context at distance $d$ from the current point of
generation.

$P(d)$ decays approximately as $1/d$ with mean slope $-1.04$ (SD $=0.15$) and no
characteristic cutoff across nearly two orders of magnitude in distance. This decay
is consistent across language families, genres, and modalities, vanishes in
random-token controls, scales with the integrity of sequential structure, and is
broadly distributed across prior sentences rather than concentrated in a small number
of anchors. The persistence function describes a continuous, scale-free temporal
structure in human language production, measurable directly from the probability
dynamics of generation and replicable across independently trained probe models
(Llama-3.1-8B, Mistral-7B).

The shape of $P(d)$ places language alongside other human behaviors and neural
processes that exhibit scale-free temporal structure, and constrains the class of
mechanisms capable of producing such structure. In particular, the absence of any
characteristic scale across the measured range --- combined with the graded
contributions of prior content distributed across the sequence --- is difficult to
reconcile with architectural accounts of memory that posit a sharp division between
maintained and retrieved material with distinguishable temporal dynamics. The pattern
is consistent instead with mechanisms in which prior content continues to shape
ongoing generation through the evolving probability structure of the sequence itself.

\section{Results}

\subsection{A persistence function for natural language}

We define the persistence function $P(d)$ as the predictive influence of prior
context at distance $d$ from the current point of generation, estimated as the
reduction in target perplexity produced by revealing prior context at that distance,
with shuffled-token baselines controlling for the contribution of token identity
alone (Methods). $P(d)$ provides a distance-resolved measurement of how preceding
language constrains the probability structure of upcoming tokens.

Across all corpora, $P(d)$ is positive at distances well beyond any plausible bound
on working-memory capacity. At distances exceeding 1000 tokens, prior context
continues to reduce target perplexity by a measurable margin relative to shuffled
controls. Long-range contextual influence is therefore not confined to a recency
window but extends as a graded function across the full sequences we examined.

\subsection{Heavy-tailed decay across languages, genres, and modalities}

The functional form of $P(d)$ is well approximated by a heavy-tailed decay. Figure~1
plots $P(d)$ against distance on log--log axes for each of the ten corpora. Across
languages, genres, and modalities, $P(d)$ decays smoothly across nearly two orders of
magnitude with a mean slope of $-1.04$ (SD $=0.15$) --- close to a $1/d$ power law
--- and a median goodness-of-fit of $r^2 = 0.96$. Fits are computed on $d \geq 10$
tokens, the range over which both ordered- and shuffled-context bases are stable;
very-short-range intervals are reported in Methods. Power-law fits substantially
outperform exponential alternatives, with $\Delta$AIC values ranging from $+12$ to
$+26$ across the majority of corpora. Two corpora, Japanese literary prose and
Buckeye spoken English, are better fit by stretched-exponential functions, but their
decay remains heavy-tailed and extends well beyond the range over which any bounded
buffer could plausibly contribute.

The cross-corpus consistency of the decay regime is notable. The sample spans
Germanic, Romance, Turkic, Japonic, Uralic, and Slavic language families, as well as
fiction and literary prose, news and expository prose, prepared speech transcripts,
and spontaneous speech. A narrow heavy-tailed regime across this range indicates that
the persistence function reflects a general property of natural language rather than a
property of any single language, register, or modality.

The qualitative pattern that defines $P(d)$ --- a positive contribution of ordered
context at long range together with a near-zero contribution of shuffled context ---
extends beyond the six families covered by the long-range run. In a complementary
cross-language analysis evaluating shorter context lengths on both probe models, the
long-range sign pattern is preserved across three additional language families:
Afro-Asiatic, Koreanic, and Sino-Tibetan (Supplementary Information). The headline
slope of $P(d)$ is therefore obtained on a six-family set within a broader nine-family
cross-language replication of the qualitative pattern that defines the persistence
function.

\subsection{The persistence function reflects sequential structure, not properties of the probe}

A central concern in any LLM-based measurement is whether the observed pattern
reflects properties of natural language or properties of the probe model. We address
this directly with synthetic-sequence controls. We constructed sequences with
token-frequency statistics matched to natural text but with sequential structure
removed. Raw perplexity over these synthetic sequences exhibits apparent power-law
structure with a slope comparable in magnitude to that observed on natural text
(Llama-3.1: $\approx -1.27$ on uniform random vocabulary; Mistral-7B: $\approx -1.34$),
replicating the general tendency for language-model perplexity to fall as context
length increases even when genuine long-range dependence is absent.

The order-specific persistence function $P(d)$, however, is near zero across all
distances on the synthetic sequences. The shuffled-token baseline that defines $P(d)$
absorbs the perplexity reduction attributable to token statistics alone, isolating the
contribution of sequential organization. The long-range, heavy-tailed $P(d)$ observed
for natural language therefore depends on properties of the linguistic sequence
itself, not on the probe. We replicated the persistence signature with two
independently trained probe models, Llama-3.1-8B and Mistral-7B, obtaining the same
qualitative pattern and comparable decay ranges on shared corpora, along with near-zero
$P(d)$ on synthetic controls. The persistence function is therefore not an idiosyncratic
property of the model used to estimate it.

\subsection{Order and content make distinguishable contributions to persistence}

To characterize the structure of $P(d)$ further, we compared ordered context to two
structure-disrupting manipulations. Sentence shuffling preserves within-sentence syntax
and lexical content while disrupting discourse-level ordering. Token shuffling
eliminates ordered structure within the contextual span.

Sentence shuffling reduces $P(d)$ at short distances, especially below approximately 30
tokens, but leaves much of the long-range tail intact (Figure~2). Token shuffling
collapses $P(d)$ toward zero across all distances. The persistence function therefore
decomposes into two distinguishable contributions: a chaining-specific component
concentrated at short range (slope $\approx -1.26$, approaching zero by $d > 100$
tokens) and a content-driven component that extends across the long-range tail (slope
$\approx -0.78$). Discourse-level chaining shapes the short-range structure of $P(d)$;
content-driven integration carries its long-range tail.

\subsection{Persistence is broadly distributed across prior sentences}

We next asked whether persistence is concentrated in a small number of privileged prior
sentences or broadly distributed across the prior context. We addressed this with a
sentence-level ablation: for each target, we removed one prior sentence at a time and
measured the resulting change in prediction.

The majority of prior sentences contribute positively to prediction at all distances
examined. The proportion of contributing sentences declines gradually with distance,
from approximately 90\% at short range to 55--60\% beyond 500 tokens, with no sharp
transition or threshold (Figure~3). The magnitude of contribution is concentrated at
recent positions (top-5 sentences account for 30--51\% of total positive influence;
Gini coefficient 0.78--0.90), but this concentration is fully explained by recency:
variation in mean influence by absolute position within the prior context is accounted
for by distance to the target, with no residual effect of structural position (e.g.,
document opening). Persistence is therefore a distributed field of influence across the
prior sequence, with magnitude graded smoothly by recency rather than localized to a
small number of anchoring positions.

\section{Discussion}

The persistence function shows a continuous, distance-dependent influence that extends
well beyond the span of explicit maintenance, decays without a characteristic scale, and
is distributed across the prior sequence rather than localized to a small number of
contributing positions. This pattern is difficult to reconcile with architectural
accounts that posit a sharp division between recently-maintained material and
remotely-retrieved material with distinguishable temporal dynamics: a sharp division
would predict a corresponding inflection in $P(d)$, and we observe none across nearly two
orders of magnitude in distance. The distributed influence pattern is similarly difficult
to reconcile with retrieval-based accounts in which long-range coherence is supported by
access to a small number of privileged anchoring positions. The data are consistent
instead with mechanisms in which prior content continues to shape ongoing generation
through the evolving probability structure of the sequence --- a graded, distance-dependent
field of influence rather than a partition between maintained and retrieved content.

The mean exponent we report ($-1.04$, SD $=0.15$) sits in a neighborhood that has been
independently identified in two adjacent traditions in memory research.
\citet{anderson1991reflections}, examining the recurrence statistics of information demand
in natural environments, reported a power law with exponent of approximately $-0.77$.
\citet{kahana2002age}, examining associative decay in free recall via lag-conditional
response probability, reported a forward-decay exponent of approximately $-0.82$ for
younger adults. Each of these measurements was made by entirely different methods on
entirely different data, and neither involves a language model. The convergence is striking
but does not by itself establish a common underlying mechanism: it motivates direct
cross-paradigm investigation of whether comparable scale-free decay forms across production,
environmental statistics, and recall reflect a domain-general dynamic of biological cognition
or distinct mechanisms producing similar exponents.

This empirical picture aligns with a longer trajectory in cognitive psychology away from
architecturally bounded views of memory and toward continuous and graded characterizations.
Cowan's focus-of-attention model \citep{cowan2001magical} recasts working memory as a graded
availability of items rather than a fixed-capacity slot store; Ericsson and Kintsch's
long-term working memory framework \citep{ericsson1995longterm} argues that skilled
performers make routine use of long-term-memory structures during ongoing task performance,
blurring the maintenance-retrieval boundary; and continuous-resource accounts in visual
working memory \citep{ma2014changing} replace discrete-item architectures with graded
distributions of representational precision. The persistence function provides a direct
quantitative measurement of an analogous continuity in language production, observed at the
scale of token-level sequential dependence and across naturalistic discourse rather than
experimental task sequences. The present result is best read as an empirical capstone to
this longer trajectory, not a refutation of it.

We do not claim that these results establish the mechanism by which prior content shapes
language production in the brain. The persistence function is a measurement of the form of
contextual influence, not of the cognitive processes that implement it. We note, however,
that artificial autoregressive systems trained on human text exhibit the same scale-free
shape under the same measurements, supporting the interpretation that $P(d)$ reflects
properties of well-formed sequential generation rather than properties specific to human
cognition. Whether the persistence function generalizes to non-linguistic sequential
behaviors with extended temporal coherence --- music improvisation, motor sequences, gesture,
drawing --- and whether the convergence with memory-retrieval exponent ranges reflects a
shared dynamic or merely shared statistical constraints, are questions the present results
motivate but do not answer.

\section{Methods}

\subsection{Corpora}

The long-range persistence function was computed on ten corpora: English fiction (Project
Gutenberg), English news prose, English spontaneous speech (Buckeye Corpus of Conversational
Speech), TED transcripts in English, German, French, Turkish, and Russian, and literary prose
in Japanese and Finnish. Together these cover six language families (Germanic, Romance, Turkic,
Slavic, Japonic, Uralic) and four production modes (written fiction, news prose, prepared
speech, spontaneous speech). The complementary cross-language qualitative replication of the
sign-flip pattern (Afro-Asiatic, Koreanic, Sino-Tibetan) was performed on Wikipedia and
TED-transcript cells from a broader corpus inventory at shorter context lengths. Full corpus
details, sources, sample sizes, and preprocessing steps are provided in Supplementary
Information.

\subsection{Probe models}

Two autoregressive language models served as probes: Llama-3.1-8B (base;
\texttt{Meta-Llama-3.1-8B}) and Mistral-7B-v0.1. Neither was fine-tuned for this study. The
probes were used only to compute conditional probability estimates over fixed token sequences;
the language models themselves were not the object of study. Headline results were obtained on
Llama-3.1-8B; the qualitative pattern was replicated on Mistral-7B for shared corpora as a
probe-independence check.

\subsection{Persistence function}

For each document, a single target region of 30 tokens was selected at the midpoint (50\%) of
document length. Context preceding the target was revealed in incrementally larger spans at
log-spaced context lengths $c \in \{0, 1, 2, 4, 8, 16, 32, 64, 128, 256, 512, 1024\}$. At each
context length, target perplexity was computed under intact context (the actual preceding $c$
tokens, original order) and shuffled context (the same $c$ tokens, random order). Per-distance
marginals were obtained as the per-token change in perplexity across successive context lengths,
$m_d = (\mathrm{ppl}_{c_{i-1}} - \mathrm{ppl}_{c_i}) / (c_i - c_{i-1})$, with distance indexed by
the geometric mean $d = \sqrt{c_{i-1} c_i}$ of the interval. The persistence function was defined
as $P(d) = m_d^{\text{intact}} - m_d^{\text{shuffled}}$, isolating the contribution of sequential
structure from probe-internal token-identity calibration. Marginals were aggregated across
documents within each corpus and binned by distance. Only documents with at least
$1024 + 30 + 50$ tokens were included, ensuring a full-length context window and target for every
measurement.

\subsection{Functional-form fitting}

The per-token order-specific gap was fit on log--log axes to four functional forms: power law
($y = \alpha\, d^{\beta}$), exponential ($y = \alpha\, e^{-\beta d}$), stretched exponential
($y = \alpha\, e^{-\beta d^{\gamma}}$), and quadratic-in-log. Models were compared by Akaike (AIC)
and Bayesian (BIC) information criteria, and functional-form preference was determined per corpus.
Power-law fits were computed on positive bins with $d \geq 10$ tokens, the first interval at which
both ordered and shuffled bases exceed eight tokens.

\subsection{Synthetic-sequence controls}

To distinguish probe-internal calibration dynamics from genuine sequential dependence, two
synthetic-sequence controls were constructed: uniform-vocabulary sequences (tokens drawn uniformly
at random from the model vocabulary) and frequency-matched sequences (tokens drawn with empirical
token-frequency probabilities matched to the natural corpora). The same persistence-function
protocol was applied to both.

\subsection{Sentence-level decomposition}

The contribution of discourse-level ordering versus within-sentence content was assessed by
computing $P(d)$ under three context conditions: ordered (intact), sentence-shuffled (sentences
permuted while within-sentence token order preserved), and token-shuffled. The chaining-specific
component was estimated as the difference between ordered and sentence-shuffled persistence; the
content-driven component as the difference between sentence-shuffled and token-shuffled. Buckeye
was excluded from this decomposition because spontaneous speech does not have unambiguous sentence
boundaries.

\subsection{Sentence-level ablation}

For documents in two corpora (English fiction and English TED transcripts, 20 documents per corpus,
1024-token prior context), each prior sentence was individually removed and the resulting change in
target perplexity was computed (leave-one-sentence-out). Per-sentence influence was characterized by
sign (positive $=$ removing the sentence increased target perplexity) and magnitude. The proportion
of contributing sentences was computed in distance bins, and concentration metrics (Gini coefficient
over per-sentence magnitudes; share of total positive influence in the top-5 sentences) were computed
per document and aggregated. To test whether magnitude variation across prior-sentence positions
reflects a privileged structural role over and above distance from the target, the full
leave-one-sentence-out influence matrix was computed on a 15-document subset of each ablation corpus;
variation in influence magnitude was accounted for by distance, with no residual effect of structural
position once distance was controlled.

\section*{Data and code availability}

Analysis code and corpus manifests are available in the project repository. Per-corpus persistence
curves and the figure-generation pipeline are provided as supplementary material.

\begin{figure}[t]
\centering
\includegraphics[width=\textwidth]{figures/fig1_persistence.pdf}
\caption{\textbf{Cross-corpus persistence functions.} Per-token order-specific contextual influence
$P(d)$ as a function of distance $d$, log--log axes. Each panel is one of ten corpora spanning six
language families (Germanic, Romance, Turkic, Slavic, Japonic, Uralic) and four production modes. All
cells exhibit heavy-tailed decay across nearly two orders of magnitude in distance, with mean power-law
slope $-1.04$ (SD $=0.15$) --- close to a $1/d$ law. Fits computed on $d \geq 10$ tokens. Cells marked
$\dagger$ (Japanese literary prose; Buckeye spontaneous English speech) are best fit by
stretched-exponential functions but remain heavy-tailed.}
\end{figure}

\begin{figure}[t]
\centering
\includegraphics[width=0.85\textwidth]{figures/fig2_sentence_shuffle.pdf}
\caption{\textbf{Order and content make distinguishable contributions to persistence.} $P(d)$ under
context conditions that isolate discourse-level ordering from within-sentence content. The intact and
sentence-shuffled curves coincide at long range and diverge only below ${\sim}30$ tokens; both lie far
above the token-shuffled baseline at all distances. The persistence function decomposes into a
chaining-specific component (slope $\approx -1.26$, approaching zero by $d > 100$ tokens) and a
content-driven component (slope $\approx -0.78$) that carries the long-range tail.}
\end{figure}

\begin{figure}[t]
\centering
\includegraphics[width=0.85\textwidth]{figures/fig3_distributed_influence.pdf}
\caption{\textbf{Persistence is broadly distributed across prior sentences.} Proportion of prior
sentences whose removal degrades target prediction as a function of distance to the target. The
proportion declines smoothly from ${\sim}90\%$ within the first 32 tokens to ${\sim}55$--$60\%$ beyond
500 tokens, with no sharp transition. The top-5 sentences account for 30--51\% of total positive
influence (Gini 0.78--0.90), but this concentration reflects recency, not a privileged structural role
for any specific position.}
\end{figure}

\bibliographystyle{plainnat}
\bibliography{references}

\end{document}
